% =======================================================
% 4. SPECIFICA DEI TEST
% =======================================================
\section{Specifica dei Test}
Questa sezione descrive i test pianificati per accertare che il sistema rispetti i requisiti individuati. Lo stato del test è indicato come \textbf{I} (Implementato), \textbf{NI} (Non Implementato), o \textbf{S} (Superato).

\subsection{Test di unità}
Verificano che le singole unità logiche del codice sorgente funzionino correttamente in modo isolato.

\begin{table}[H]
	\centering
	\renewcommand{\arraystretch}{1.3}
	\begin{tabularx}{\textwidth}{|c|X|c|}
		\hline
		\rowcolor{primarycolor!10}
		\textbf{ID} & \textbf{Descrizione} & \textbf{Stato} \\
		\hline
		TU1 & {[Descrizione del test di unità]} & NI \\
		\hline
	\end{tabularx}
	\caption{Test di Unità}
\end{table}

\subsection{Test di integrazione}
Garantiscono che la comunicazione tra i moduli (es. Frontend e Backend) avvenga in modo coerente.

\begin{table}[H]
	\centering
	\renewcommand{\arraystretch}{1.3}
	\begin{tabularx}{\textwidth}{|c|X|c|}
		\hline
		\rowcolor{primarycolor!10}
		\textbf{ID} & \textbf{Descrizione} & \textbf{Stato} \\
		\hline
		TI1 & {[Descrizione del test di integrazione]} & NI \\
		\hline
	\end{tabularx}
	\caption{Test di Integrazione}
\end{table}

\subsection{Test di sistema}
Confermano che il software assembrato soddisfi i requisiti funzionali dichiarati nell'AdR.

\begin{table}[H]
	\centering
	\renewcommand{\arraystretch}{1.3}
	\begin{tabularx}{\textwidth}{|c|X|c|c|}
		\hline
		\rowcolor{primarycolor!10}
		\textbf{ID} & \textbf{Descrizione} & \textbf{Req. Rif.} & \textbf{Stato} \\
		\hline
		TS1 & {[Descrizione del test e requisito collegato]} & RF1 & NI \\
		\hline
	\end{tabularx}
	\caption{Test di Sistema}
\end{table}

\subsection{Test di accettazione}
Eseguiti per confermare formalmente alla proponente che il software è pronto per il rilascio.

\begin{table}[H]
	\centering
	\renewcommand{\arraystretch}{1.3}
	\begin{tabularx}{\textwidth}{|c|X|c|}
		\hline
		\rowcolor{primarycolor!10}
		\textbf{ID} & \textbf{Descrizione} & \textbf{Stato} \\
		\hline
		TA1 & {[Descrizione test collaudo finale]} & NI \\
		\hline
	\end{tabularx}
	\caption{Test di Accettazione}
\end{table}

\subsection{Tracciamento test-requisiti}
Tutti i requisiti obbligatori devono essere tracciati con almeno un Test di Sistema (TS) e un Test di Accettazione (TA).